% Section 4.1: 整体框架与实验流程

第三章已得到可执行单测判定、带思维步骤与算法标签的统一语料。本章在此基础上搭建方法层对照，要点如下：
\begin{itemize}
    \item \textbf{语料与任务口径}：统一字段与可执行单测判定，保证训练/评测可比；
    \item \textbf{共同基座}：CodeGeeX4-ALL-9B（\texttt{THUDM/codegeex4-all-9b}），权重落盘于 \texttt{models/CodeGeeX4-ALL-9B/}；
    \item \textbf{析因因子}：训练侧「是否及如何注入 LoRA」与推理侧「是否采用两阶段 CoT 及 few-shot 规模」两条独立维度；
    \item \textbf{分析目标}：构造覆盖主效应与交互效应的实验矩阵，为第五章 pass@k 归因提供完整对照基础。
\end{itemize}

\paragraph{（1）两因素析因设计}
对照维度严格区分为两个正交因子，避免训练侧与推理侧的效应彼此混淆：
\begin{itemize}
    \item \textbf{训练侧（LoRA 适配器类型）} 共四档——\texttt{Base}（不加 LoRA 的原始基座）、\texttt{LoRA\_code}（仅监督 \texttt{solution} 的纯代码 LoRA）、\texttt{LoRA\_cot\_zs}（监督 \texttt{thought\_step + solution} 且训练 prompt 不含示例）、\texttt{LoRA\_cot\_fs(k)}（同样监督 CoT 格式但训练 prompt 含 $k$ 条 few-shot 示例）；
    \item \textbf{推理侧（Prompt 结构）} 共三档——\texttt{Direct}（单阶段，题目直接映射到代码）、\texttt{CoT-ZS}（两阶段，先生成推理再生成代码，零示例）、\texttt{CoT-FS(k)}（在 CoT-ZS 基础上前置 $k$ 条按强标签检索得到的示例）。
\end{itemize}

\paragraph{（1a）训练侧：四类权重配置下，单条样本在 \texttt{apply\_chat\_template} 之前的 messages 骨架}
约定与实现对应关系如下：
\begin{itemize}
    \item \textbf{占位符}：\texttt{\{func\}}、\texttt{\{params\}}、\texttt{\{question\}} 来自 \texttt{test\_info[0]} 与题面；\texttt{\{thought\}}、\texttt{\{code\}} 对应 \texttt{thought\_step}、\texttt{solution}；
    \item \textbf{脚本对应}：\texttt{LoRA\_code} $\leftrightarrow$ \texttt{train/train.py}；\texttt{LoRA\_cot\_zs} / \texttt{LoRA\_cot\_fs} $\leftrightarrow$ \texttt{train/train\_cot.py}；
    \item \textbf{模板再包装}：下列为语义块，最终字节串仍由基座 \texttt{apply\_chat\_template} 封装。
\end{itemize}

\begin{description}
  \item[\texttt{Base}]
    \begin{itemize}
        \item 不更新 LoRA，无训练 messages；评测加载 \texttt{models/CodeGeeX4-ALL-9B/}。
    \end{itemize}
  \item[\texttt{LoRA\_code}] 单条样本为\textbf{一轮}目标对话（无 few-shot）。
    \begin{itemize}
        \item \textbf{System}：与 Direct 推理相同——\texttt{You are an expert Python programmer. Write correct Python code to solve the given problem.}
        \item \textbf{User}（末尾禁止解释，仅要代码）：
\begin{quote}\ttfamily\footnotesize
Implement the function `\{func\}' that takes \{params\} to solve:\\[0.3em]
\{question\}\\[0.3em]
Only output the Python code, no explanation.
\end{quote}
        \item \textbf{Assistant}：可执行 Python 正文，\textbf{无} Markdown 围栏，与 \texttt{solution} 一致；
        \item \textbf{Loss}：仅最后一轮 Assistant 段。
    \end{itemize}
  \item[\texttt{LoRA\_cot\_zs}] 单条样本仍为\textbf{一轮}目标对话，$k{=}0$。
    \begin{itemize}
        \item \textbf{System}（强调先推理后编码）：
\begin{quote}\ttfamily\footnotesize
You are an expert Python programmer. Before writing code, think through\\
the problem step by step. Format your response as numbered reasoning\\
steps followed by the Python solution.
\end{quote}
        \item \textbf{User}（CoT 版题面指令）：
\begin{quote}\ttfamily\footnotesize
Implement the function `\{func\}' that takes \{params\} to solve:\\[0.3em]
\{question\}\\[0.3em]
Think step by step, then output the Python code.
\end{quote}
        \item \textbf{Assistant}：\texttt{thought\_step} + 空行 + \verb|```python| \texttt{solution} \verb|```|（无 \texttt{thought} 时退化为仅代码块）；
        \item \textbf{Loss}：仅最后一轮 Assistant 段。
    \end{itemize}
  \item[\texttt{LoRA\_cot\_fs($k$)}] 在 \texttt{LoRA\_cot\_zs} 的 System 下拼接 \textbf{$k$ 组示范}后再接当前题。
    \begin{itemize}
        \item \textbf{示范结构}：每组为 \textbf{User(示例题，CoT 版指令)} $\rightarrow$ \textbf{Assistant(示例 \texttt{thought+code}，格式同 \texttt{LoRA\_cot\_zs})}，与 \texttt{build\_few\_shot\_messages} 一致；
        \item \textbf{当前题}：再接一条 User（同上 CoT 版指令）+ 一条 Assistant（当前题 \texttt{thought+code}）；
        \item \textbf{Loss}：\textbf{仅}当前题 Assistant；前 $k$ 组示范在 labels 中全部为 \texttt{-100}。
    \end{itemize}
\end{description}

\paragraph{（1b）推理侧：三种 Prompt 结构下一次评测请求长什么样}
\begin{itemize}
    \item \textbf{\texttt{Direct}}（\texttt{eval\_base\_passk.py} / \texttt{eval\_lora\_passk.py}）
    \begin{itemize}
        \item \textbf{调用形态}：单阶段、一次 \texttt{generate}；
        \item \textbf{前缀}：\textbf{System} + \textbf{User} 与（1a）中 \texttt{LoRA\_code} 训练前缀\textbf{逐字同构}（含 \texttt{Only output the Python code...}）；
        \item \textbf{后处理}：从生成文本抽取围栏内代码，送单测；
        \item \textbf{对应组别}：H1 / H4 / H7 / H9，差异仅在是否挂载 LoRA 及挂载哪一种适配器。
    \end{itemize}
    \item \textbf{\texttt{CoT-ZS}} 与 \textbf{\texttt{CoT-FS($k$)}}（\texttt{evaluation/eval\_lora\_cot\_passk.py}）
    \begin{itemize}
        \item \textbf{调用形态}：两阶段、两次 \texttt{generate}；
        \item \textbf{LoRA}：可选；\texttt{--lora-path NONE} 时为基座上的 CoT（H2/H3）；
        \item \textbf{Few-shot}：$k{>}0$ 时，两阶段使用\textbf{同一组} $k$ 条示例；评分与训练侧一致，抽样种子 \texttt{2026 + problem\_idx*1000 + sample\_idx}；
        \item \textbf{Stage-1}（只生成自然语言推理）
        \begin{itemize}
            \item \textbf{System}：
\begin{quote}\ttfamily\footnotesize
You are an expert algorithm teacher. Generate step-by-step reasoning.
\end{quote}
            \item \textbf{若 $k{>}0$}：依次追加 $k$ 组 \textbf{User}（示例的 CoT 版题指令）$\rightarrow$ \textbf{Assistant}（示例 \texttt{thought\_step} + fenced \texttt{solution}）；
            \item \textbf{目标题}：再追加一条 \textbf{User}（CoT 版题指令，与训练 \texttt{LoRA\_cot} 相同）；
            \item \textbf{清洗}：若输出中出现 \verb|```python|，截断其后，避免推理段「偷跑」代码。
        \end{itemize}
        \item \textbf{Stage-2}（在固定推理条件下生成代码）
        \begin{itemize}
            \item \textbf{System}：
\begin{quote}\ttfamily\footnotesize
You are an expert Python programmer. Based on the reasoning, write the code.
\end{quote}
            \item \textbf{若 $k{>}0$}：再次回放与 Stage-1 \textbf{相同顺序}的 $k$ 组 User $\rightarrow$ Assistant（示例全文）；
            \item \textbf{目标综合 User}：写入截断后的 Stage-1 文本，并重复函数签名与题面，末尾限定 fenced Python：
\begin{quote}\ttfamily\footnotesize
Reasoning:\\[0.2em]
\{truncated\_stage1\_text\}\\[0.5em]
Now implement the function `\{func\}' that takes \{params\} to solve:\\[0.3em]
\{question\}\\[0.3em]
Write only the Python code in \verb|```python| block.
\end{quote}
        \end{itemize}
    \end{itemize}
    \item \textbf{与 H1--H10 的读法对照}
    \begin{itemize}
        \item H2 / H5 / H8：两阶段、$k{=}0$（Stage-1/2 均无示例回放）；
        \item H3 / H6 / H10：两阶段、$k{>}0$（两阶段各多 $k$ 组示范轮次）；
        \item H7 / H9：推理侧为 Direct，但训练侧 messages 仍按（1a）中 CoT 或 CoT+few-shot 构造，与（1）节因子划分一致。
    \end{itemize}
\end{itemize}

\paragraph{（2）合法性约束与有效组合}
\begin{itemize}
    \item \textbf{组合规模}：训练侧 4 档 $\times$ 推理侧 3 档，共 12 种理论组合；
    \item \textbf{核心约束}：若训练侧与推理侧\textbf{均为} CoT，则两侧 few-shot 条数 $k$ 必须\textbf{严格相等}，以保持 train/test 上下文分布一致，避免「训练 $k{=}0$、推理 $k{>}0$」或反之带来的漂移。
\end{itemize}
据此排除以下两组非法组合：
\begin{itemize}
    \item \texttt{LoRA\_cot\_zs} + \texttt{CoT-FS(k>0)}：训练 $k{=}0$ 与推理 $k{>}0$ 不一致；
    \item \texttt{LoRA\_cot\_fs(k)} + \texttt{CoT-ZS}：训练 $k{>}0$ 与推理 $k{=}0$ 不一致。
\end{itemize}
剩余 \textbf{10 组有效实验} H1--H10，构成本文全部对照的源头。

\paragraph{（3）十组实验矩阵}
表~\ref{tab:ch4-matrix} 汇总 H1--H10，表中口径如下：
\begin{itemize}
    \item \textbf{$k$ 取值}：与脚本 \texttt{FS\_K} 一致，本文固定为 $2$；
    \item \textbf{LoRA 指针}：\texttt{train/} 下纯文本文件，记录最近一次训练落盘的适配器目录，供评测读取；
    \item \textbf{列「评测脚本」}：与仓库 \texttt{evaluation/} 中入口及关键参数一一对应。
\end{itemize}

\begin{table}[!htb]
  \centering
  \caption{H1--H10 实验矩阵：训练侧 $\times$ 推理侧与评测入口}
  \label{tab:ch4-matrix}
  \footnotesize
  \setlength{\tabcolsep}{3pt}
  \begin{tabular}{clllll}
    \toprule
    编号 & 训练侧 & 推理侧 & LoRA 指针 & 评测脚本 \\
    \midrule
    H1  & Base            & Direct      & ---                                    & \texttt{eval\_base\_passk.py} \\
    H2  & Base            & CoT-ZS      & ---                                    & \texttt{eval\_lora\_cot\_passk.py --lora-path NONE --few-shot-k 0} \\
    H3  & Base            & CoT-FS($k$) & ---                                    & \texttt{eval\_lora\_cot\_passk.py --lora-path NONE --few-shot-k k} \\
    H4  & LoRA\_code      & Direct      & \texttt{latest\_lora\_adapter.txt}       & \texttt{eval\_lora\_passk.py} \\
    H5  & LoRA\_code      & CoT-ZS      & \texttt{latest\_lora\_adapter.txt}       & \texttt{eval\_lora\_cot\_passk.py --lora-path <code> --few-shot-k 0} \\
    H6  & LoRA\_code      & CoT-FS($k$) & \texttt{latest\_lora\_adapter.txt}       & \texttt{eval\_lora\_cot\_passk.py --lora-path <code> --few-shot-k k} \\
    H7  & LoRA\_cot\_zs   & Direct      & \texttt{latest\_lora\_cot\_zs\_adapter.txt} & \texttt{eval\_lora\_passk.py --lora-path <cot\_zs>} \\
    H8  & LoRA\_cot\_zs   & CoT-ZS      & \texttt{latest\_lora\_cot\_zs\_adapter.txt} & \texttt{eval\_lora\_cot\_passk.py --few-shot-k 0} \\
    H9  & LoRA\_cot\_fs($k$) & Direct   & \texttt{latest\_lora\_cot\_fs\{k\}\_adapter.txt} & \texttt{eval\_lora\_passk.py --lora-path <cot\_fs>} \\
    H10 & LoRA\_cot\_fs($k$) & CoT-FS($k$) & \texttt{latest\_lora\_cot\_fs\{k\}\_adapter.txt} & \texttt{eval\_lora\_cot\_passk.py --few-shot-k k} \\
    \bottomrule
  \end{tabular}
\end{table}

\paragraph{（4）整体实验流水线}
\begin{itemize}
    \item \textbf{目标}：10 组结果可复现、相互可比；
    \item \textbf{训练阶段}：\texttt{train.py} 产出 \texttt{LoRA\_code}；\texttt{train\_cot.py}（$k{=}0$ / $k{>}0$）分别产出 \texttt{LoRA\_cot\_zs}、\texttt{LoRA\_cot\_fs($k$)}；各脚本将最新适配器目录写入独立指针文件；
    \item \textbf{评测阶段}：按开关调用 \texttt{eval\_base\_passk.py}、\texttt{eval\_lora\_passk.py}、\texttt{eval\_lora\_cot\_passk.py}，读取对应指针，结果统一写入 \texttt{evaluation/results/eval\_*\_<suffix>.json}；
    \item \textbf{示意图}：数据流见图~\ref{fig:ch4-pipeline}。
\end{itemize}

\begin{figure}[!htb]
  \centering
  \begin{tikzpicture}[
    node distance=5mm and 6mm,
    box/.style={draw, rounded corners=2pt, align=center, minimum height=8mm, minimum width=24mm, font=\footnotesize},
    lbl/.style={font=\scriptsize\itshape, align=center},
    arr/.style={-{Stealth[length=2mm]}, thick},
  ]
    % row 1: training
    \node[box] (d) {\texttt{KodCode\_train.jsonl}};
    \node[box, right=of d] (t1) {\texttt{train.py}};
    \node[box, right=of t1] (t2) {\texttt{train\_cot.py}\\ \scriptsize($k{=}0$)};
    \node[box, right=of t2] (t3) {\texttt{train\_cot.py}\\ \scriptsize($k$)};

    \node[box, below=9mm of t1] (a1) {\texttt{LoRA\_code}};
    \node[box, below=9mm of t2] (a2) {\texttt{LoRA\_cot\_zs}};
    \node[box, below=9mm of t3] (a3) {\texttt{LoRA\_cot\_fs($k$)}};

    \draw[arr] (d) -- (t1);
    \draw[arr] (d.east) to [bend left=10] (t2.west);
    \draw[arr] (d.east) to [bend left=20] (t3.west);
    \draw[arr] (t1) -- (a1);
    \draw[arr] (t2) -- (a2);
    \draw[arr] (t3) -- (a3);

    % row 2: evaluation
    \node[box, below=11mm of a1, minimum width=34mm] (e1) {\texttt{eval\_base\_passk.py}\\ \scriptsize H1};
    \node[box, right=of e1, minimum width=34mm] (e2) {\texttt{eval\_lora\_passk.py}\\ \scriptsize H4 / H7 / H9};
    \node[box, right=of e2, minimum width=34mm] (e3) {\texttt{eval\_lora\_cot\_passk.py}\\ \scriptsize H2 / H3 / H5 / H6 / H8 / H10};

    \draw[arr] (a1.south) to [bend right=10] (e2.north);
    \draw[arr] (a2.south) to [bend right=10] (e2.north);
    \draw[arr] (a3.south) to [bend right=10] (e2.north);
    \draw[arr] (a1.south) to [bend right=20] (e3.north);
    \draw[arr] (a2.south) to [bend right=10] (e3.north);
    \draw[arr] (a3.south) to [bend left=10] (e3.north);

    \node[box, below=9mm of e2, minimum width=50mm] (r) {\texttt{evaluation/results/eval\_*\_<suffix>.json}};
    \draw[arr] (e1.south) to [bend right=10] (r.west);
    \draw[arr] (e2) -- (r);
    \draw[arr] (e3.south) to [bend left=10] (r.east);
  \end{tikzpicture}
  \caption{第四章整体流水线：训练阶段产出三份适配器，评测阶段按 10 组配置统一落盘 pass@k 结果}
  \label{fig:ch4-pipeline}
\end{figure}

\paragraph{（5）结果命名与主要对照口径}
\begin{itemize}
    \item \textbf{文件命名}：\texttt{eval\_<mode>\_passk\_<hN\_desc>.json}，后缀与表~\ref{tab:ch4-matrix} 末列一致；
    \item \textbf{元数据}：CoT 评测 JSON 另存 \texttt{mode}、\texttt{use\_lora}、\texttt{lora\_path}、\texttt{few\_shot\_k} 等字段，便于回溯；
    \item \textbf{主效应}：以 H1 为共同基线，分别通过 $\{$H4, H7, H9$\}$ 抽取训练侧主效应，通过 $\{$H2, H3$\}$ 抽取推理侧主效应；
    \item \textbf{交互效应}：由 $(H8-H7)-(H2-H1)$ 与 $(H10-H9)-(H3-H1)$ 两条差分差分量，判定 LoRA 与 CoT 是否正交可加；
    \item \textbf{Few-shot 一致性}：对比 $H3-H2$ 与 $H6-H5$ 的符号与量级，判断 few-shot 边际增益是否依赖训练侧。
\end{itemize}
本章后续安排：
\begin{itemize}
    \item \textbf{4.2--4.4 节}：依次给出 Base、LoRA、CoT 三条分支的实现细节；
    \item \textbf{超参口径}：除表~\ref{tab:ch4-matrix} 所列「训练侧 $\times$ 推理侧」取值外，其余训练与评测超参全局固定，不随组别改动。
\end{itemize}